\documentclass[final,12p]{article}
\usepackage{graphicx}
\usepackage[letterpaper,margin=1.0in]{geometry}
\usepackage{amssymb}
\usepackage{changepage}
\usepackage{float}
\usepackage{hyperref}
\usepackage{url}
\usepackage{afterpage}
\usepackage{natbib}
\usepackage{setspace}
\usepackage{fancyhdr}
\pagestyle{fancy}
\fancyhf{}
\usepackage[utf8]{inputenc} 
\usepackage{lastpage}


\def\Student{Julio C\'{e}sar Borb\'{o}n Fragoso}
\def\Title{ANTEPROYECTO}
\def\Prog{Doctorado en Ciencias (F\'{i}sica) }
\def\Dept{Departamento de Investigac\'{i}on en Fis\'{i}ca}
\def\Institution{Universidad de Sonora}
\def\Director{}
\def\ProjectTitle{}
\def\ResearchLine{E}

%%header and footer
\lhead{\Student\ / \Prog}
\rhead{\Title}
\rfoot{Page \thepage \hspace{1pt} of \pageref{LastPage}}


%%%%%%%comands
\newcommand{\SubItem}[1]{ {\setlength\itemindent{15pt} \item[-] #1} }
  

%%%%%%%%%%%%%%%%%%%%%%%%%%%%%%%%%%%%%
\begin{document}
\onehalfspacing

%%%%Title Page
\begin{titlepage}
\centering
\hspace{0pt}
\vfill
{\scshape\Large \Title \par}
%project revised 2023-2
  
  \vspace{2cm}
  {
    TITLE:\par
    {\bf \large \ProjectTitle \par}
  }
       
  \vspace{0.5cm}
  {
    RESEARCH LINE: \par
    \ResearchLine \par
  }
        
  \vspace{4cm}
  {\underline{\hspace{8cm}}\par}
  {\bf \scshape \Student \par}
  {Student\par}

  \vspace{1cm}
  {\underline{\hspace{8cm}}\par}
  {\scshape \Director \par}
  {Director\par}

  \vspace{1cm}
  {\bf \Prog \par}
  {\Dept \par}
  {\Institution \par}

  \vspace{4cm}
  {\today}

\hspace{0pt}
\vfill

\end{titlepage}


%%%%% white page for print out
\shipout\null


%%%% Abstract Page
\newpage
\hspace{2pt}
\vfill

  \begin{center}
    {\Large \ProjectTitle \par}
    \vspace{1cm}
    {\itshape\textbf{Abstract}\par}
  \end{center}
  
  \vspace{0.5cm}
   
This research 

  \hspace{2pt}
\vfill

%%%%Empieza el cuerpo
\newpage
\section{Marco teorico}

Portales oscuros: 

$\mathcal{L} = - \frac{1}{2} B^{\mu \nu} F^{'}_{\mu \nu}$ Foton oscuro $A^{'}_{\nu}$ cineticamente mezclado con hipercarga.

$\mathcal{L} = (AS + \lambda S^2) H^t H$ Escalar oscuro S acoplado al Higgs.

\section{Propuesta}

Estudiar el sector oscuro con observables que pueden ser medidos en detectores de ondas gravitacionales o aceleradores de particulas. 

\section{Objetivo General}

\section{Hipotesis}

\section{Metodologia}

\section{Objetivos especificos}

\section{Resultados esperados}

\section{Calendario de actividades}

\section{Habilidades a desarrollar por el estudiante}

\section{Mobilidad}

\end{document}
